% Ad Atticum 6.2 (ad-atticum-06-02) --- English translation
% Translated by Claude (Anthropic) from the Latin (Perseus).
% Style guide: STYLE.md.

\ciceroLetterOpener{CICERO ATTICO salutem}

\ciceroSection{1} When your freedman Philogenes had come to me at Laodicea to pay his respects, and said he would sail to you at once, I gave him this letter, in which I am answering those of yours that I received by Brutus's courier. And I will reply first to the last page of yours, which gave me great vexation --- I mean what Cincius wrote to you about Statius's talk. The most vexing thing in it is that Statius says this plan is approved by me also. Approved? Of the matter itself I will say only this much: that for my own part I could wish for many further bonds of the closest connection between us, although the bonds of affection are already the tightest possible; so far am I from wishing any loosening of that by which we are bound together.

\ciceroSection{2} As for him, I have often had occasion to know that he speaks rather sharply about these things, and I have often soothed him when he was angry. You know this, I suppose. And in this travelling-and-soldiering life of mine I have often seen him fired up with anger, and often pacified. What he wrote to Statius I do not know. Whatever he meant to do in such a matter, he certainly ought not to have written to a freedman. It shall be my own greatest concern that nothing turn out otherwise than we wish and than is proper. Nor is it enough in a matter of this kind that each man answer for himself: the largest share of that duty belongs to the boy Cicero, or rather the young man already; which is the very thing I am always urging on him. And he does seem to me to love his mother very much, as he should, and yourself wonderfully. But the boy's nature, while great, is still many-sided; and in managing it I have business enough.

\ciceroSection{3} Since I have answered your last page with my first, I now come back to your first. That all the cities of the Peloponnese were on the sea I had believed on the authority of Dicaearchus's tables --- a writer not without merit and, by your own judgment, well-approved. On many counts, in his Trophonian narrative of Chaeron, he reproaches the Greeks with this, that they pursued nothing but the sea; and he does not except a single place in the Peloponnese. Since the authority pleased me (for he was \textit{[Greek: very thoroughly historical]} and had lived in the Peloponnese), I nevertheless wondered, and, scarcely able to believe it, consulted Dionysius. He was at first taken aback; then, since he thought no worse of Dicaearchus than you do of C. Vestorius or I do of M. Cluvius, he had no doubt that we should trust him. He held that there was a certain coastal Lepreon in Arcadia; Tenea, Aliphera, and Tritia, on the other hand, seemed to him \textit{[Greek: new foundations]}, and he confirmed this by \textit{[Greek: the Catalogue of Ships]}, where no mention is made of them. I therefore took over that passage from Dicaearchus word for word. That they are called Phliasians I knew, and that is how you should have it; that, at any rate, is what we have. But at first \textit{[Greek: analogy]} had misled me --- \textit{[Greek: Phlious, Opous, Sipous]}, since one says \textit{[Greek: Opountioi, Sipountioi]}. But this we straight away corrected.

\ciceroSection{4} That you take pleasure in our moderation and restraint I see well. You would do so the more if you were on the spot. Indeed, in this assize-circuit which I held from the Ides of February at Laodicea to the Kalends of May, of every district except Cilicia, we have achieved certain marvellous things. So many cities have been freed of all their debts, many have been greatly relieved; all of them, in the use of their own laws and judgments, have come back to life, having obtained \textit{[Greek: autonomy]}. I have given them the means to free or relieve themselves of debt in these two ways: first, that absolutely no expense has been incurred in my province --- and when I say none, I am not speaking \textit{[Greek: hyperbolically]} --- none, I tell you, not so much as a farthing.

\ciceroSection{5} It is incredible how much, by this alone, the cities have surfaced. Then another thing was added. Astonishing thefts had been committed in the cities by their own Greeks, by their own magistrates. I myself examined those who had held magistracies in the last ten years. They confessed openly. So, with no disgrace put on them, they paid the money back to their peoples out of their own shoulders. And the peoples, without a groan, paid back to the publicans the arrears of the previous lustrum as well --- to the very men to whom in this same lustrum they had paid nothing. And so with the publicans I am in their good books. ``Grateful men,'' you will say. We have felt it. As for the rest of the jurisdiction, it has been neither unskilled nor lenient, but combined with a remarkable accessibility; and access to me has been least of all provincial in character --- nothing routed through a chamberlain; before daybreak I have been walking about the house as I used to in my candidate days. These things are welcome and great, and not yet burdensome to me, thanks to that old soldiering.

\ciceroSection{6} On the Nones of May I had it in mind to set out for Cilicia. Having spent the month of June there (and would that it were in peace! for a great war hangs over us from the Parthians), I meant to put July into the return journey. For my annual term of duty runs out on the third day before the Kalends of Sextilis. I am, however, in great hope that no extension of my time will be granted. I have had the city register down to the Nones of March, from which I gather that, through the firmness of our friend Curio, anything will be done rather than the business of the provinces. So, as I hope, I shall see you soon.

\ciceroSection{7} I come to your Brutus --- or rather ours, as you prefer to have it. I have indeed done everything I could either to settle in my own province or to attempt in the kingdom. With the king I have therefore dealt in every way, and I deal with him daily by letter. For I had him in person with me for three or four days, in troubled circumstances, from which I freed him. But I did not cease, both then when he was at hand and afterwards in the most frequent letters, to ask and press my own case; to advise and exhort him for his own. I have gained much, but how much I do not plainly know, since I am far away. As for the Salaminians (these I could constrain), I brought them to the point of being willing to pay the whole sum to Scaptius, but with interest reckoned at twelve percent per annum from the date of the most recent bond, and not running continuously but renewed every year. The money was being counted out; Scaptius would not have it. You who say that Brutus is keen to lose something? He had forty-eight percent in his bond. It could not be done; nor, if it could, could I bear it. I hear that Scaptius is altogether sorry now. For as to the decree of the Senate he alleged --- that judgment be given on the bond --- it was passed for this reason: that the Salaminians had taken the money in defiance of the Lex Gabinia. But the law of Aulus forbade judgment to be given on money taken in such a way. Therefore the Senate decreed that judgment should be given on this particular bond. As things now stand, that bond has the same legal standing as any other, and nothing privileged.

\ciceroSection{8} These things, done by me in due order, I expect to satisfy Brutus; whether to satisfy you I do not know; Cato I shall certainly satisfy. But I come back now to yourself. Tell me, Atticus --- you, the praiser of my integrity and elegance --- did you dare, with your mouth (as Ennius says), to ask me to give cavalry-troops to Scaptius for the collection of the money? If you were with me here --- you who write that you are sometimes vexed at not being so --- would you let me do this if I wished? You say, ``not more than fifty.'' Spartacus had fewer than that at first. What mischief would such men not have done in so tender a little island? Would not have done? On the contrary, what did they not do before my arrival? They kept the Salaminian Senate shut up in the council-house so many days that some of them died of starvation. For Appius's prefect was Scaptius, and he had squadrons from Appius. Do you, then --- you whose face, by Hercules, is always before my eyes whenever I think of any duty or any honour --- do you, I say, ask me to make Scaptius a prefect? We had laid it down on other occasions that no businessman should hold the office, and Brutus had accepted this. Is he to have squadrons? Why squadrons rather than cohorts? In point of expense Scaptius is turning into a real spendthrift.

\ciceroSection{9} ``The leading men want it,'' you say. I know; for they came to me all the way at Ephesus and tearfully told me of the cavalry's crimes and their own miseries. And so I at once gave orders that the cavalry should be out of Cyprus by a fixed day; and for that, then, and for the rest, the Salaminians lifted us to the skies in their decrees. But why cavalry now? For the Salaminians are paying; unless perhaps we wish to achieve, by force of arms, that they should reckon their interest at forty-eight percent. And shall I dare ever to read or to touch those books of yours which you so loudly praise, if I do such a thing? Too much, my dearest Atticus, have you in this matter loved Brutus; me, I fear, too little. And what I have here written, I have written to Brutus, as having been written to me by you.

\ciceroSection{10} Take now the rest of the news. For Appius I am doing here everything --- with honour, yet plainly with goodwill. For I do not hate the man himself, and I love Brutus; and Pompey is wonderfully urgent with me on his behalf, a man whom by Hercules I love more and more from day to day. You have heard that C. Caelius the quaestor is coming out here. I do not know what kind of man he is. But that Pammenian business of yours does not please me. I hope myself to be at Athens in the month of September. The dates of your travels I should very much like to know. The \textit{[Greek: simple-mindedness]} of Sempronius Rufus I learned from your letter from Corcyra. What more would you have? I envy the power of Vestorius. I should have liked even now to prattle on more; but it is light; the crowd presses; Philogenes is in a hurry. Farewell, then; bid Pilia and our Caecilia farewell by letter; and a greeting to you from my Cicero.
